% Imporditud: tools/import_eksperimendid.py
% Allikas: Eesti füüsikaolümpiaadi eksperimendi ülesannete
%   kogu (Rael Kalda, 2023), ülesanne Ü60, lahendus L60.
\setAuthor{EFO žürii}
\setRound{piirkonnavoor}
\setYear{2013}
\setNumber{G E1}
\setDifficulty{2}
\setTopic{varia}
\setVahendid{1- ja 20-sendised mündid, valge A4 joonistuspaber, pliiats. Joonlauda kasutada ei tohi!}
\setPunktid{8}
\setAllikas{Eksperimendi kogu Ü60, lahendus lk 59}
\setLahendusseis{toores}

\prob{Sendid}
Mitu korda erinevad 1- ja 20-eurosendiste müntide ruumalad?

\hint

\solu
Mündi ruumala saame leida valemist V = Sph (0,5 p), kus Sp = $\pi$$d_2$/4 (0,5 p).

Seega ruumalade suhte saame leida, kui leiame mündi diameetrite suhte ning pak-

suste suhte V20 d220 h20 $V_1$ d21 $h_1$ = (0,5 p).

Diameetrite suhte leidmiseks vaatame, mitu korda mahuvad mündid A4 paberi ühe külje peale (1,5 p). Paberi pikema külje peale mahub 18,3 1-sendist ja 13,3 20-sendist, seega d20/$d_1$ = 18,3/13,3 = 1,38 (1 p). Diameetri suhte leidmiseks võib veeretada münte paberi peal ning vaadata mitu täisringi moodustab üks paberi külg. (Diameetrite suhte leidmine iga sarnase täpsusega meetodiga (2,5 p). Müntide paksuste leidmiseks voldime paberi kokku ning vaatame mitu kihti paberit on võrdne mü(ndi paksus)ega (1 p). Joonistuspaberiga 120 g/cm$^2$ saame, et 1-sendine münt on 12 lehte paks ning 20-sendine münt on 15 lehte paks. Seega h20/$h_1$ = 15/12 = 1,25. (1 p) Müntide ruumalade suhe on seega V20/$V_1$ = 1,382 $\cdot$ 1,25 = 2,38. Tulemus vahemikus 2,3 - 2,5 annab 2 p, vahemikus 2,1 - 2,7 annab 1 p.

Märkus: Täpne müntide ruumalade suhe on 2,4, kuid antud meetodiga tuleb ebatäpsus sisse müntide paksuste suhte leidmisel.
\probend
